\section{Methodology}\label{sec:method}

\subsection{Recursive Bayesian estimation}
The navigation state at time step $k$ is
\begin{equation}
  \xv_k = \begin{bmatrix} \pv_k \\ \bias_k \end{bmatrix}
       = \begin{bmatrix} p_{x,k} & p_{y,k} & b_{x,k} & b_{y,k} \end{bmatrix}\trans,
\end{equation}
the horizontal position $\pv_k$ (easting and northing in UTM zone~44N, metres) and the
velocity bias $\bias_k$ of the inertial navigation system (INS, m/s). Given the INS
outputs $\uu_{1:k}$ and the terrain-height measurements $z_{1:k}$, the goal is the filtering
density $p(\xv_k \mid z_{1:k})$, computed by the Bayes recursion~\cite{arulampalam2002,sarkka2023,thrun2005}
\begin{align}
  p(\xv_k \mid z_{1:k-1}) &= \int p(\xv_k \mid \xv_{k-1}, \uu_k)\,
                            p(\xv_{k-1} \mid z_{1:k-1})\, d\xv_{k-1}
                            && \text{(prediction)} \label{eq:chapman}\\
  p(\xv_k \mid z_{1:k}) &= \frac{p(z_k \mid \xv_k)\, p(\xv_k \mid z_{1:k-1})}
                               {p(z_k \mid z_{1:k-1})}
                            && \text{(update)}. \label{eq:bayes}
\end{align}
Every estimator in this project implements \eqref{eq:chapman}--\eqref{eq:bayes} under a
different approximation: the particle filter with a weighted sample, the EKF with a Gaussian
obtained by linearising the terrain, and dead reckoning by skipping the update.

\subsection{System and sensor models}\label{sec:models}

\paragraph{Inertial navigation (motion model).}
The INS delivers a velocity $\uu_k = \vect{v}_k + \bias_k + \vect{n}_k$, i.e.\ the true
velocity corrupted by a slowly wandering bias and white noise. Integrating
$\vect{v}_k = \uu_k - \bias_k - \vect{n}_k$ over a step $\Delta t$ and modelling the bias as a
random walk gives the linear motion model
\begin{equation}
  \xv_k = F\xv_{k-1} + G\uu_k + \vect{w}_k,\qquad
  F = \begin{bmatrix} I_2 & -\Delta t\, I_2 \\ 0 & I_2 \end{bmatrix},\quad
  G = \begin{bmatrix} \Delta t\, I_2 \\ 0 \end{bmatrix}.
  \label{eq:motion}
\end{equation}
\Cref{eq:motion} is exact for the simulated world (\cref{sec:sim}) when $\vect{w}_k$ is the
correlated noise derived in \cref{sec:pcrb}. The filters use the simpler diagonal
$Q = \mathrm{diag}(\sigma_p^2 I_2, \sigma_b^2 I_2)$ with $\sigma_p = 2$\,m per step, much
larger than the true position noise ($\approx 0.05$\,m per step): a deliberate tuning choice
that keeps the particle cloud diverse and absorbs unmodelled effects. This matters when the
filters' error is compared with the bound in \cref{sec:pcrb}.

\paragraph{Altimeters (measurement model).}
A barometric altimeter measures the altitude above mean sea level, $a_k + e_{b,k}$, and a radar
altimeter the height above the surface below, $a_k - h(\pv_k) + e_{r,k}$. Their difference
cancels the unknown altitude $a_k$:
\begin{equation}
  z_k = \underbrace{(a_k + e_{b,k})}_{\text{baro}} - \underbrace{(a_k - h(\pv_k) + e_{r,k})}_{\text{radar}}
      = h(\pv_k) + v_k,\qquad v_k = e_{b,k} - e_{r,k} \sim \Normal(0, \sigma_z^2),
  \label{eq:meas}
\end{equation}
with $\sigma_z^2 = \sigma_{\mathrm{baro}}^2 + \sigma_{\mathrm{radar}}^2 = 2^2 + 3^2$\,m$^2$
($\sigma_z \approx 3.6$\,m). Both noise levels, their whiteness and their independence of
range are \emph{assumptions} of this study, not values from a sensor datasheet. The terrain
height $h(\pv)$ is read from the digital elevation model (DEM) by bilinear interpolation, which
is exact for a plane and continuous everywhere. The likelihood $p(z_k \mid \xv_k) =
\Normal(z_k; h(\pv_k), \sigma_z^2)$ is a nonlinear function of position and, because many places
share the same height, it is \emph{multimodal}: a single measurement constrains the position
only to a contour line of the terrain.

\subsection{The bootstrap (SIR) particle filter}\label{sec:pf}
A particle filter represents the posterior by $N$ weighted samples
$\{\xv_k^{(i)}, w_k^{(i)}\}_{i=1}^N$,
$p(\xv_k \mid z_{1:k}) \approx \sum_i w_k^{(i)} \delta(\xv_k - \xv_k^{(i)})$
\cite{gordon1993,arulampalam2002,doucet2011}. With samples drawn from a proposal
$q(\xv_k \mid \xv_{k-1}^{(i)}, z_k)$, sequential importance sampling gives the weight recursion
\begin{equation}
  w_k^{(i)} \propto w_{k-1}^{(i)}\,
  \frac{p(z_k \mid \xv_k^{(i)})\, p(\xv_k^{(i)} \mid \xv_{k-1}^{(i)})}
       {q(\xv_k^{(i)} \mid \xv_{k-1}^{(i)}, z_k)}.
\end{equation}
The \emph{bootstrap} filter of Gordon, Salmond and Smith~\cite{gordon1993} uses the motion
model as the proposal, $q = p(\xv_k \mid \xv_{k-1})$, so the weight update reduces to the
likelihood:
\begin{equation}
  \xv_k^{(i)} \sim p(\xv_k \mid \xv_{k-1}^{(i)}, \uu_k),\qquad
  w_k^{(i)} \propto w_{k-1}^{(i)}\, p(z_k \mid \xv_k^{(i)}).
  \label{eq:sir}
\end{equation}
Each particle carries its own bias hypothesis $\bias^{(i)}$ and integrates the INS velocity
corrected by it; particles with a wrong bias drift away from the terrain that matches the
measurements and lose weight, so the bias is estimated without being measured.

\paragraph{Log-domain weights.} A residual of 50\,m at $\sigma_z = 3.6$\,m multiplies a weight by
$e^{-96}$; a few such factors underflow double precision. The filter therefore stores
$\ell^{(i)} = \log w^{(i)}$ and normalises with the log-sum-exp identity
\begin{equation}
  \ell^{(i)} \leftarrow \ell^{(i)} - \log\sum_j e^{\ell^{(j)}}, \qquad
  \log\sum_j e^{\ell^{(j)}} = m + \log\sum_j e^{\ell^{(j)} - m},\; m = \max_j \ell^{(j)}.
\end{equation}
The normaliser $\log\sum_j w_{k-1}^{(j)} p(z_k\mid\xv_k^{(j)})$ is the Monte Carlo estimate of
the log evidence $\log p(z_k \mid z_{1:k-1})$ in \eqref{eq:bayes}; it is logged for the
divergence monitor of \cref{sec:monitor}.

\paragraph{Estimate.} The minimum mean-square-error estimate and its covariance are
\begin{equation}
  \hat\xv_k = \sum_i w_k^{(i)} \xv_k^{(i)},\qquad
  P_k = \sum_i w_k^{(i)} (\xv_k^{(i)} - \hat\xv_k)(\xv_k^{(i)} - \hat\xv_k)\trans .
\end{equation}
When the posterior is multimodal the mean can lie between modes, where there are no
particles; this is expected behaviour of the MMSE estimate, not an error.

\paragraph{Degeneracy and resampling.} Repeated multiplication by likelihoods concentrates the
weight on a few particles. The effective sample size
\begin{equation}
  \Neff = \frac{1}{\sum_i (w_k^{(i)})^2} \in [1, N]
\end{equation}
measures this; when $\Neff < \tau N$ (default $\tau = 0.5$) the filter \emph{resamples}: it
draws $N$ new particles with replacement, particle $i$ receiving $N_i$ copies with
$\Exp[N_i] = N w^{(i)}$, and resets the weights to $1/N$. Four unbiased schemes are
implemented~\cite{hol2006}; they differ in the variance of $N_i$, which is extra noise added to
the estimate:
\begin{itemize}
  \item \textbf{multinomial}: $N$ independent draws from $\mathrm{Categorical}(w)$, so
        $N_i \sim \mathrm{Bin}(N, w^{(i)})$ and $\mathrm{Var}(N_i) = N w^{(i)}(1-w^{(i)})$;
  \item \textbf{stratified}: one uniform draw in each stratum, $u_j = (j - 1 + U_j)/N$;
  \item \textbf{systematic}: one uniform $U$ shared by all strata, $u_j = (j - 1 + U)/N$;
        each particle receives $\lfloor N w^{(i)} \rfloor$ or $\lceil N w^{(i)} \rceil$
        copies;
  \item \textbf{residual}: $\lfloor N w^{(i)} \rfloor$ copies deterministically, the remaining
        $R = N - \sum_i \lfloor N w^{(i)} \rfloor$ by multinomial sampling on the residual
        weights.
\end{itemize}
All four are implemented by inverting the cumulative weights at sorted pointers $u_j$.

\paragraph{Roughening.} Resampling duplicates particles, and with small process noise the cloud
can collapse to few distinct values (sample impoverishment). Roughening~\cite{gordon1993} adds
independent jitter after each resampling, with per-component standard deviation
$\sigma_j = K E_j N^{-1/d}$, where $E_j$ is the range of component $j$ over the particles, $d = 4$
the state dimension and $K$ a tuning constant.

\subsection{Extended Kalman filter baseline}\label{sec:ekf}
The EKF approximates the posterior by a Gaussian $\Normal(\hat\xv_k, P_k)$. Because the motion model
\eqref{eq:motion} is linear its prediction is exact; only the measurement is linearised, at the
predicted position, through the terrain slope:
\begin{align}
  \hat\xv_k^- &= F\hat\xv_{k-1} + G\uu_k, & P_k^- &= F P_{k-1} F\trans + Q,\\
  H_k &= \begin{bmatrix} \partial h/\partial p_x & \partial h/\partial p_y & 0 & 0\end{bmatrix}_{\hat\pv_k^-},
  & S_k &= H_k P_k^- H_k\trans + \sigma_z^2,\\
  K_k &= P_k^- H_k\trans S_k^{-1}, & \hat\xv_k &= \hat\xv_k^- + K_k\,(z_k - h(\hat\pv_k^-)),\\
  P_k &= (I - K_k H_k) P_k^- (I - K_k H_k)\trans + K_k \sigma_z^2 K_k\trans. &&
\end{align}
The Joseph form of the covariance update keeps $P_k$ symmetric positive semi-definite. The
slope comes from central differences of the DEM, bilinearly interpolated. The weakness is
structural: $H_k$ describes the terrain over one 30\,m cell, while the position uncertainty may
be hundreds of metres. Over that range the terrain is nothing like a plane, so the EKF explains a
large innovation with the wrong local slope and can move confidently in the wrong direction.
It also cannot represent several candidate positions at once.

\subsection{Heavy-tailed likelihood for altimeter outliers}\label{sec:studentt}
A radar altimeter occasionally returns a gross error (multipath, the edge of a cliff, a
spurious return). Under the Gaussian likelihood a single 100\,m outlier multiplies the weight
of every particle near the truth by about $e^{-386}$ (since $\tfrac12(100/3.6)^2 \approx 386$),
wiping them out. Replacing the Gaussian
by a Student-$t$ density with scale $\sigma_z$ and $\nu$ degrees of freedom,
\begin{equation}
  \log p(z \mid \xv) = \log\Gamma\!\left(\tfrac{\nu+1}{2}\right) - \log\Gamma\!\left(\tfrac{\nu}{2}\right)
  - \tfrac12 \log(\nu\pi\sigma_z^2)
  - \tfrac{\nu+1}{2}\log\!\left(1 + \frac{(z - h(\pv))^2}{\nu\sigma_z^2}\right),
  \label{eq:studentt}
\end{equation}
makes the log-likelihood grow only logarithmically with the residual, so one outlier costs
a few units instead of hundreds and the particles near the truth survive. For small residuals
the core of \eqref{eq:studentt} is close to the Gaussian, so little accuracy is lost when there
are no outliers. The constants are kept so that the log evidence of the two models can be
compared; $\nu = 4$ is used.

\subsection{Divergence monitor and recovery}\label{sec:monitor}
A filter that has locked onto the wrong patch of terrain (``false lock'') is confidently
wrong: its cloud is tight, yet the terrain under it does not match the measurements. This is
detectable without ground truth through the normalised innovation squared (NIS) of the
predictive distribution,
\begin{equation}
  \mathrm{NIS}_k = \frac{(z_k - \hat z_k)^2}{S_k},\quad
  \hat z_k = \sum_i w_{k-1}^{(i)} h(\pv_k^{(i)}),\quad
  S_k = \sum_i w_{k-1}^{(i)} \bigl(h(\pv_k^{(i)}) - \hat z_k\bigr)^2 + \sigma_z^2 .
\end{equation}
For a filter consistent with the data, $\mathrm{NIS}_k$ is approximately $\chi^2_1$, so the sum
over a window of $L$ steps is approximately $\chi^2_L$; the monitor raises an alarm when the
windowed mean exceeds $\chi^2_L(1 - p_{\mathrm{FA}})/L$. During acquisition the cloud is wide,
$S_k$ is large and the NIS stays small, so the monitor is naturally quiet there. On an alarm the
belief is replaced by the mixture
\begin{equation}
  (1-\alpha)\, p(\xv_k \mid z_{1:k}) + \alpha\, p_{\mathrm{DR}}(\xv_k),
\end{equation}
represented by $(1-\alpha)N$ particles resampled from the current cloud and $\alpha N$ particles
drawn from the dead-reckoning belief $p_{\mathrm{DR}}$ --- the prior propagated through the
motion model without any terrain update. Whatever the terrain did, $p_{\mathrm{DR}}$ still
contains the true state at its stated confidence, so the re-spread particles give the filter a
chance to find the correct mode again, while the kept particles preserve the current hypothesis
in case the alarm was false.

\subsection{Rao-Blackwellised (marginalised) particle filter}\label{sec:rbpf}
The bias enters the model linearly and the measurement does not depend on it. Conditioned on a
position history $\pv_{0:k}$ the bias is therefore Gaussian, and Rao-Blackwellisation computes it
exactly with a Kalman filter instead of sampling it~\cite{schon2005}: particles carry only the
position, each with the Kalman mean $\hat\bias^{(i)}$ of its bias. The Kalman \emph{covariance}
$P^b$ depends only on the model matrices, which are the same for every particle, so all
particles share one covariance and one Riccati recursion. Marginalising the bias out of
\eqref{eq:motion} gives the position proposal, and the sampled increment
$\vect{\eta}^{(i)} = \pv_k^{(i)} - \pv_{k-1}^{(i)} - \Delta t\,\uu_k = -\Delta t\,\bias_{k-1} + \vect{w}_p$
is a linear measurement of the bias:
\begin{align}
  \pv_k^{(i)} &\sim \Normal\bigl(\pv_{k-1}^{(i)} + \Delta t(\uu_k - \hat\bias^{(i)}),\; S\bigr),
  & S &= \Delta t^2 P^b + \sigma_p^2 I_2,\\
  L &= -\Delta t\, P^b S^{-1}, & \hat\bias^{(i)} &\leftarrow \hat\bias^{(i)} + L\bigl(\vect\eta^{(i)} + \Delta t\,\hat\bias^{(i)}\bigr),\\
  P^b &\leftarrow P^b - L S L\trans + \sigma_b^2 I_2, & w_k^{(i)} &\propto w_{k-1}^{(i)}\, p(z_k \mid \pv_k^{(i)}).
\end{align}
The state covariance adds the shared $P^b$ to the spread of the bias means. Each particle now
represents a whole Gaussian over the bias instead of one sample of it, which should need fewer
particles for the same accuracy.

\subsection{Posterior Cramér--Rao bound}\label{sec:pcrb}
The posterior (Bayesian) Cramér--Rao bound gives a lower bound on the mean-square error of
\emph{any} estimator, $\Exp[(\hat\xv_k - \xv_k)(\hat\xv_k - \xv_k)\trans] \succeq J_k^{-1}$, where $J_k$
is the Bayesian information matrix. Tichavsk\'y, Muravchik and Nehorai~\cite{tichavsky1998}
derived a recursion for $J_k$; for linear dynamics with additive Gaussian noise and the
measurement \eqref{eq:meas} it becomes
\begin{equation}
  J_k = \bigl(Q + F J_{k-1}^{-1} F\trans\bigr)^{-1} + \Exp\!\left[H_k\trans R^{-1} H_k\right],\qquad
  J_0 = P_0^{-1},\qquad H_k = \begin{bmatrix} \nabla h(\pv_k)\trans & 0 & 0 \end{bmatrix},
  \label{eq:pcrb}
\end{equation}
where the Jacobian is evaluated at the \emph{true} state and the expectation is over the
distribution of the true state. The first term is the Kalman prediction of the information; the
second is the information the altimeter adds, $\nabla h \nabla h\trans / R$: steep terrain is
informative, and only along the slope direction. Bergman applied this bound to terrain
navigation~\cite{bergman1999thesis}. For a linear-Gaussian model \eqref{eq:pcrb} is exactly the
Kalman filter covariance recursion; the implementation is tested against it.

\paragraph{True noise.} The bound must use the noise of the simulated world, not the filters'
tuned $Q$. From $\uu_k = \vect v_k + \bias_k + \vect n_k$ and
$\bias_k = \bias_{k-1} + \Delta\bias_k$,
$\pv_k = \pv_{k-1} + \Delta t(\uu_k - \bias_{k-1}) - \Delta t(\Delta\bias_k + \vect n_k)$, so the
process noise $\vect w_k = [-\Delta t(\Delta\bias_k + \vect n_k);\ \Delta\bias_k]$ has, per axis,
\begin{equation}
  Q_{\mathrm{true}} = \begin{bmatrix} \Delta t^2 (q_b + q_n) & -\Delta t\, q_b \\
                                      -\Delta t\, q_b & q_b \end{bmatrix}\otimes I_2,
  \qquad q_b = \sigma_{\mathrm{rw}}^2 \Delta t,\; q_n = \sigma_n^2,
\end{equation}
and $R = \sigma_z^2$. A unit test checks $Q_{\mathrm{true}}$ against the empirical covariance of
the simulator's noise.

\paragraph{Which expectation.} Two versions are computed. (i) The \emph{ensemble bound}: the
expectation in \eqref{eq:pcrb} is the average over the true trajectories of all Monte Carlo
runs. Flights are placed uniformly at random, so given the initial guess the true start is
distributed exactly as the filters' prior $\Normal(\hat\pv_0, \sigma_0^2 I)$, and this bound applies to
the ensemble RMSE the experiments report. (ii) The \emph{per-flight bound}: $H_k$ along one fixed
flight path, compared with the tracking error of runs that fly that path. Because the bound only
involves the slope at the true position, it ignores ambiguity --- other places with the same
height profile. It describes a filter that already knows which hill it is over, so it should be
tight in the tracking phase and loose during acquisition.

\subsection{Terrain data}\label{sec:data}
The terrain is the Copernicus GLO-30 digital elevation model~\cite{copernicus2022}, a
1-arc-second (\(\approx 30\)\,m) global product, downloaded through the OpenTopography Global
DEM API~\cite{opentopo} and reprojected with bilinear resampling to a 30\,m grid in UTM zone~44N
(EPSG:32644), which covers all of Sri Lanka. The product handbook states an absolute vertical
accuracy better than 4\,m and a relative vertical accuracy better than 2\,m (slopes up to
20\,\%) or 4\,m (steeper slopes), both as 90\,\% linear errors~\cite{copernicus2022}. It is a
\emph{surface} model: it includes buildings and vegetation, so over forest a radar altimeter
that sees the ground (or part of the canopy) measures something different from the map. Study
regions were chosen by measured terrain statistics with a rule fixed in advance
(\cref{sec:regions}). The terrain statistics used throughout are:
the elevation range and standard deviation; the slope $|\nabla h|$ in percent; the
\emph{roughness}, the standard deviation of elevation inside a moving $1\times1$\,km window,
averaged over the region; and the \emph{Fisher information density}
$\Exp[|\nabla h|^2]/\sigma_z^2$, the average position information one altimeter measurement adds
(\cref{eq:pcrb}).

\subsection{Evaluation methodology}\label{sec:eval}
\paragraph{Monte Carlo with common random numbers.} Every result is an average over $M$ Monte
Carlo runs. Run $i$ draws its flight (random start, heading and weave phase), INS noise,
altimeter noise, outliers, map error and initial-guess direction from separate named random
streams seeded by $i$, with fixed-size draws. Every filter and every configuration therefore
sees exactly the same flights and noise: differences between estimators are paired, and a
parameter change alters nothing but what it scales.

\paragraph{Metrics.} A run has \emph{converged} if its position error stays below 50\,m
throughout at least the final 100\,s; it is a \emph{false lock} if the final error exceeds
200\,m while the filter's own $2\sigma$ position radius $2\sqrt{\lambda_{\max}(P_{\mathrm{pos}})}$
is below 50\,m (confidently wrong). A run can be neither: a late acquisition or a temporary
excursion while crossing flat ground counts as ``not converged'' even if the final error is
small. The \emph{tracking RMSE} is computed over the final 100\,s of converged runs only, which
separates tracking accuracy from acquisition. Consistency is measured by the position NEES
$\epsilon_k = \tilde\pv_k\trans P_{\mathrm{pos},k}^{-1}\tilde\pv_k$ and its average over runs (ANEES),
which for a consistent filter lies in
$[\chi^2_{2M}(0.025)/M,\ \chi^2_{2M}(0.975)/M]$~\cite{barshalom2001}.

\paragraph{Uncertainty and tests.} Rates carry Wilson score 95\,\% intervals~\cite{wilson1927};
RMSE values carry percentile-bootstrap 95\,\% intervals over runs~\cite{efron1993}. Because the
runs are paired, two filters' convergence is compared with McNemar's exact
test~\cite{mcnemar1947}, which uses only the runs where exactly one of them converged. All
parameters of the analysis (selection rules, definitions, which regions each experiment uses)
were fixed and recorded before the corresponding experiment ran.
